\section{Cota cuantitativa de la capa fija: Bockstein, clases y torsión}
\label{sec:capa-fija}

La Nota~\ref{sec:finitud-uniforme} probó la finitud del presupuesto
para un cuerpo y un corte fijos. El paso siguiente exige una estimación
que conserve sus dependencias cuando cambian esos datos. Desarrollamos
la deducción de la cota de capa fija de A93: un conteo homológico
controla valores asociados a relaciones de divisores; una comparación
con clases de ideales cancela el factor que depende de la altura;
una estimación integral de torsión proporciona el tamaño final.

Las pruebas elementales se dan completas. La realización aritmética de
los ciclos, el rango de una imagen homológica, la estabilidad integral,
los mapas de comparación con clases y Selmer y la cota de torsión
se reciben con sus enunciados precisos. El teorema final demuestra
la implicación desde esas entradas; permite localizar qué falta
reconstruir para verificar cada una de ellas.

\fuente{A93, Proposición~9.4, Lemas~9.5--9.7 y desarrollo,
pp.~67--74, ecuaciones (9.7)--(9.20). Se conserva el ejemplar
\texttt{295153115a03}. Milne~\cite{milne-ant}, Teoremas~4.3 y~8.31,
pp.~70 y~147, aporta las entradas clásicas de Minkowski y Chebotarev.
Soulé~\cite{soule-torsion}, §§4.1--4.3, pp.~778--779, delimita
el antecedente de torsión fuera de pequeños primos. Las menciones
de interfaces de A93 se atribuyen a esa fuente; este bloque no añade
lecturas ni ejecución de Lean.}

\subsection{La estimación que se quiere obtener}

Sea $L$ totalmente imaginario, $n=[L:\Q]$, $D=|\operatorname{disc}L|$
y $r_2=r_2(L)=n/2$. Para un primo $\ell$ y $e\ge1$, mantenemos
$s_{\ell,e}(L)$ y $W_\ell(L,e)$ de \eqref{eq:fu-grupos}.
El conjunto $S_\ell$ sigue constando de los lugares sobre $\ell$
y los infinitos. El objetivo cuantitativo es
\begin{equation}
 \log\#H^2(G_{L,S_\ell},\Z_\ell(e+1))
 =W_\ell(L,e)\log\ell
 \le \exp\!\left((2en)^{C_K}\log(2D)\right),
 \label{eq:cf-objetivo}
\end{equation}
con un exponente absoluto $C_K$. En esta ruta se supone
$\mu_\ell\subset L$ si $\ell$ es impar y $\mu_4\subset L$ si
$\ell=2$. Estas hipótesis pertenecen al argumento cuantitativo;
la finitud clásica de la Nota~\ref{sec:finitud-uniforme} no las necesitaba.

Escribiremos
\[
 \Gamma_N=\operatorname{GL}_N(\mathcal O_L),\qquad
 T_{N,j}(L)=\#H_j(\Gamma_N,\Z)_{\mathrm{tors}}.
\]
Siempre que usemos esos órdenes recibimos su finitud. La estimación
de su tamaño se introducirá después del conteo. El entero
$Q=\ell^\nu$ designa el módulo de los coeficientes; no es el cuerpo
$\Q$. La altura $\nu$ debe desaparecer de la cota de clases antes
de tomar el paso aritmético a Selmer.

\subsection{Cadenas y un prisma de conjugación}

Para un grupo $G$ y un anillo conmutativo $A$, el complejo bar
con coeficientes triviales tiene $C_m(G,A)$ libre sobre los símbolos
$[g_1|\cdots|g_m]$, con $C_0(G,A)=A$. Para $m\ge1$,
\begin{equation}
 \begin{aligned}
 \partial[g_1|\cdots|g_m]
 &=[g_2|\cdots|g_m]\\
 &\quad+\sum_{i=1}^{m-1}(-1)^i
 [g_1|\cdots|g_ig_{i+1}|\cdots|g_m]\\
 &\quad+(-1)^m[g_1|\cdots|g_{m-1}].
 \end{aligned}
 \label{eq:cf-bar}
\end{equation}
En grado cero la frontera es cero. Se usa el complejo no normalizado:
los símbolos con una entrada igual a la identidad no se eliminan.
Las identidades entre caras emparejan con signos opuestos los
términos de $\partial^2$, de modo que $\partial^2=0$.
Un homomorfismo de grupos actúa sobre cada entrada y conmuta con
la frontera.

\begin{lema}[Prisma integral]
\label{lem:cf-prisma}
Para $h\in G$, sea $c_h(g)=h^{-1}gh$. El operador
\begin{equation}
 P_h[g_1|\cdots|g_m]
 =\sum_{i=0}^{m}(-1)^i
 [g_1|\cdots|g_i|h|c_h(g_{i+1})|\cdots|c_h(g_m)]
 \label{eq:cf-prisma}
\end{equation}
satisface $\partial P_h+P_h\partial=(c_h)_*-1$.
Por tanto la conjugación induce la identidad en homología.
\end{lema}

\begin{proof}
Denotemos por $H_i$ el sumando de \eqref{eq:cf-prisma} sin signo.
Una cara estrictamente anterior a la inserción coincide con la
inserción correspondiente después de aplicar esa cara a la cadena
original. Una cara posterior cumple la misma identidad porque
$c_h$ respeta productos. Esos términos de $\partial P_h$ se
cancelan con los de $P_h\partial$.

Las caras que unen $h$ a una entrada vecina se cancelan entre
$H_i$ y $H_{i+1}$: la identidad necesaria es
$h\,c_h(g)=gh$. Quedan la primera cara de $H_0$, que aplica
$c_h$ a todas las entradas, y la última de $H_m$, cuyo signo es
$(-1)^m(-1)^{m+1}=-1$. En grado cero,
$\partial[h]=0$ y ambos extremos actúan igual en $A$.
Esto prueba la fórmula en todos los grados. Todas las operaciones
son compatibles con homomorfismos que transporten también $h$.
\end{proof}

En grado uno se ve el signo directamente:
$P_h[g]=[h|h^{-1}gh]-[g|h]$ tiene frontera
$[h^{-1}gh]-[g]$. Esta identidad explica cómo unir prismas.
Todavía hace falta realizarlos en los estabilizadores de los
retículos aritméticos pertinentes; esa realización será una entrada
distinta.

\subsection{Bockstein y levantamiento de un ciclo reducido}

\begin{lema}[Cambio de coeficientes]
\label{lem:cf-bockstein}
Sea $C_\bullet$ un complejo no negativo de grupos abelianos libres,
$Q\ge1$ y $m\ge1$. Existe una sucesión exacta natural
\begin{equation}
 0\longrightarrow H_m(C_\bullet)/QH_m(C_\bullet)
 \longrightarrow H_m(C_\bullet/QC_\bullet)
 \xrightarrow{\ \beta\ }H_{m-1}(C_\bullet)[Q]
 \longrightarrow0.
 \label{eq:cf-bockstein}
\end{equation}
Si el Bockstein de un ciclo reducido es cero, ese mismo ciclo
tiene un levantamiento integral que es ciclo.
\end{lema}

\begin{proof}
Levantemos un ciclo reducido $\bar z$ a $z\in C_m$.
Existe un único $y\in C_{m-1}$ con $\partial z=Qy$.
Como $Q\partial y=0$ y $C_{m-2}$ no tiene torsión,
$y$ es ciclo; para $m=1$ tomamos $C_{-1}=0$.
Definimos $\beta[\bar z]=[y]$.
Cambiar el levantamiento por $z+Qv$ cambia $y$ por $y+\partial v$.
Cambiar $\bar z$ por una frontera añadida permite usar
$z+\partial u+Qv$ y tampoco altera la clase. La aplicación está
bien definida y $Q[y]=[\partial z]=0$.

Si $[y]$ está anulado por $Q$, existen representantes con
$\partial z=Qy$. Su reducción es ciclo y tiene Bockstein $[y]$;
esto da la sobreyectividad. Si $[y]=0$, escribimos
$y=\partial w$ y ponemos $x=z-Qw$. Entonces $\partial x=0$
y $x$ se reduce exactamente a $\bar z$. No es necesario alterar
el ciclo reducido por una frontera.

Finalmente, si un ciclo integral $x$ se reduce a una frontera,
podemos escribir $x=\partial u+Qv$. Aplicar $\partial$ y usar
que $C_{m-1}$ no tiene torsión muestra que $v$ es ciclo.
Así $[x]=Q[v]$. La inclusión contraria es inmediata y demuestra
la inyectividad del primer mapa de \eqref{eq:cf-bockstein}.
Las construcciones conmutan con mapas de complejos.
\end{proof}

No se supuso generación finita de los grupos de cadenas ni de
sus homologías. Para $Q=1$ todos los términos de
\eqref{eq:cf-bockstein} son cero. La obstrucción que aparece a la
derecha tiene grado $m-1$; será la primera de las dos contribuciones
de torsión al conteo.

\subsection{Contar imágenes por rango y contar por fibras}

\begin{proposicion}[Dos conteos finitos]
\label{prop:cf-rango-fibras}
Sea $H$ abeliano con torsión finita de orden $T$, sea $J\subset H$
de rango a lo sumo $r$, y sea $f:H\to B$ un homomorfismo con
$QB=0$. Entonces
\begin{equation}
 \#f(J)\le Q^rT.
 \label{eq:cf-rango}
\end{equation}
Además, si $\beta:Y\to A$ y $g:Y\to B$ son homomorfismos,
$A$ es finito y cada subconjunto finito de $g(\ker\beta)$ tiene
a lo sumo $C$ elementos, entonces
\begin{equation}
 \#g(Y)\le C\,\#\beta(Y)\le C\,\#A.
 \label{eq:cf-fibras}
\end{equation}
No se exige que $H,J,Y$ o $B$ sean finitos o finitamente generados.
\end{proposicion}

\begin{proof}
Para un conjunto finito $E\subset f(J)$, elijamos una preimagen
en $J$ de cada elemento y llamemos $J_0$ al subgrupo que generan.
Ahora sí $J_0$ es finitamente generado. Su descomposición es
$\Z^{r_0}\oplus T_0$, con $r_0\le r$ y
$T_0\subset H_{\mathrm{tors}}$. Como $f$ anula $QJ_0$,
\[
 \#E\le\#(J_0/QJ_0)
 =Q^{r_0}\#(T_0/QT_0)\le Q^rT.
\]
Si $f(J)$ tuviera más elementos, un subconjunto de cardinal
$Q^rT+1$ contradiría esa desigualdad.

Para el segundo conteo basta tomar un conjunto finito de valores
y elegir preimágenes en $Y$. Separamos esas preimágenes por
su valor bajo $\beta$. Al restar un representante en cada fibra,
sus diferencias quedan en $\ker\beta$; los valores de $g$ en
esa fibra forman una traslación de un subconjunto de
$g(\ker\beta)$. Hay a lo sumo $C$ valores por fibra y
$\#\beta(Y)$ fibras. El mismo argumento de subconjuntos finitos
da \eqref{eq:cf-fibras} para toda la imagen.
\end{proof}

La torsión usada en \eqref{eq:cf-rango} es la del grupo ambiente
$H$. La imagen de un grupo libre puede ser un grupo de torsión:
$\Z\to\Z/Q\Z$ ya lo muestra. El teorema de estructura se
aplicó antes de tomar la imagen.

\subsection{El operador que selecciona la imagen de rango pequeño}

Fijemos un grado $k+1$ y $Q=\ell^\nu$. Para cada $N$ suponemos
dada una aplicación $\Z/Q$-lineal
$\Phi_N:C_{k+1}(\Gamma_N,\Z/Q)\to B$, donde $QB=0$.
Se exige que anule fronteras, respete la estabilización en ciclos
y satisfaga
\begin{equation}
 \Phi_{dN}(\Psi_d z)=d\Phi_N(z)
 \quad\text{para ciclos }z,
 \label{eq:cf-primitividad}
\end{equation}
donde $\Psi_d$ repite una matriz en $d$ bloques diagonales.
Reducir ciclos integrales y evaluar induce
$\phi_N:H_{k+1}(\Gamma_N,\Z)\to B$.

Para $x\in H_{k+1}(\Gamma_N,\Z)$ definimos
\begin{equation}
 z_0=x,\qquad
 z_{j+1}=(\Psi_\ell)_*z_j-\ell^{j+2}\operatorname{st}_*z_j
 \in H_{k+1}(\Gamma_{\ell^{j+1}N},\Z).
 \label{eq:cf-filtrado}
\end{equation}
Ambos términos se transportan al mismo tamaño de matrices.

\begin{lema}[Valor del operador filtrado]
\label{lem:cf-filtrado}
Para $j\ge0$,
\begin{equation}
 \phi_{\ell^jN}(z_j)
 =\ell^j u_j\phi_N(x),\qquad
 u_j=\prod_{i=1}^j(1-\ell^i).
 \label{eq:cf-valor-filtrado}
\end{equation}
En particular, cada valor $\ell^k\phi_N(x)$ pertenece a la
imagen bajo $\phi_{\ell^kN}$ del subgrupo generado por los $z_k$.
\end{lema}

\begin{proof}
Para $j=0$ el producto es vacío y vale uno. La primitividad
\eqref{eq:cf-primitividad} y la estabilización muestran que
el paso siguiente multiplica el valor por
$\ell-\ell^{j+2}=\ell(1-\ell^{j+1})$.
La inducción prueba \eqref{eq:cf-valor-filtrado}.
Cada factor de $u_k$ es primo a $\ell$. Existe por ello un entero
$b$ con $bu_k\equiv1\pmod Q$. Como $QB=0$,
$\phi_{\ell^kN}(bz_k)=\ell^k\phi_N(x)$.
\end{proof}

\paragraph{Entrada de rango y estabilidad.}
A93, pp.~67--68, recibe dos afirmaciones integrales. Para $k$ par,
$k+1<r_2$ y $N\ge2(k+1)+4$, el subgrupo generado por los $z_k$
tiene rango a lo sumo $r_2$. Además, estabilizar desde tamaño
$N_0\ge2j+4$ induce un isomorfismo de homología integral en
grado $j$ hacia cualquier tamaño mayor.
La primera afirmación concierne a esa imagen filtrada, no a toda
la homología impar. La segunda preserva la torsión; una
comparación únicamente racional no bastaría.

\subsection{Por qué aparecen dos grados de torsión}

\begin{proposicion}[Conteo homológico]
\label{prop:cf-conteo-homologico}
Recibamos las dos entradas anteriores y la finitud de la torsión.
Sea $k=2e$, $N_0=4e+6$ y $2e+1<r_2$. Para una familia
$\Phi_N$ como antes, todo conjunto finito de valores
$\ell^{2e}\Phi_N(z)$, con $z$ ciclo reducido y tamaños $N$
posiblemente distintos, tiene cardinal a lo sumo
\begin{equation}
 Q^{r_2}T_{N_0,2e}(L)T_{N_0,2e+1}(L).
 \label{eq:cf-conteo-homologico}
\end{equation}
El módulo $B$ puede ser un producto infinito de módulos anulados
por $Q$.
\end{proposicion}

\begin{proof}
Estabilicemos los finitos ciclos elegidos a un mismo tamaño
$N\ge N_0$. La aplicación de Bockstein toma valores en
$H_k(\Gamma_N,\Z)[Q]$, un grupo de cardinal a lo sumo
$T_{N_0,k}(L)$ por estabilidad integral.

Dentro de una fibra, la diferencia de dos ciclos tiene Bockstein
cero. Por el Lema~\ref{lem:cf-bockstein} se levanta a un ciclo
integral $x$. La diferencia de los valores contados es
$\ell^k\phi_N([x])$. El Lema~\ref{lem:cf-filtrado} la sitúa
en la imagen del subgrupo filtrado en
$H_{k+1}(\Gamma_{\ell^kN},\Z)$.
Ese subgrupo tiene rango a lo sumo $r_2$, y la torsión del grupo
ambiente tiene orden $T_{N_0,k+1}(L)$ por estabilidad.
La Proposición~\ref{prop:cf-rango-fibras} acota los valores por
fibra por $Q^{r_2}T_{N_0,k+1}(L)$.

Multiplicar por el número de fibras da
\eqref{eq:cf-conteo-homologico}. Solo se contó un conjunto finito
de valores; nunca se usó la finitud del módulo de llegada.
\end{proof}

El grado $2e$ aporta las fibras de Bockstein y el grado $2e+1$
aporta la torsión dentro de cada fibra. La potencia $\ell^{2e}$
proviene del operador filtrado. Son tres contribuciones diferentes.

\subsection{Caracteres equivariantes y norma ponderada}

Sea $c$ máximo tal que $\mu_{\ell^c}\subset L$, con $c\ge1$
y $c\ge2$ si $\ell=2$. Para $\nu\ge c$, consideremos una
extensión cíclica $F/L$ de grado $m_\nu=\ell^{\nu-c}$, una raíz
primitiva $\zeta$ de orden $Q=\ell^\nu$ y un generador
$\gamma$ con $\gamma(\zeta)=\zeta^{1+\ell^c}$.
Ponemos $\kappa=1+\ell^c$ y $a=\kappa^e$.
El grupo de caracteres pertinente es
\begin{equation}
 X_\nu=
 \{\varphi\in\operatorname{Hom}(\operatorname{Cl}(F),\Z/Q):
   \varphi(\gamma^{-1}\mathfrak a)=a\varphi(\mathfrak a)
   \text{ para toda clase }\mathfrak a\}.
 \label{eq:cf-caracteres}
\end{equation}
La condición obliga a factorizar por
$\operatorname{Cl}(F)/(\gamma^{-1}-a)\operatorname{Cl}(F)$;
el coeficiente $\Z/Q$ obliga además a anular $Q\operatorname{Cl}(F)$.
Para cualquier grupo abeliano finito $C$,
\[
 \#\Hom(C,\Z/Q)=\#(C/QC).
\]
En efecto, los factores primarios distintos de $\ell$ no contribuyen,
y un factor $\Z/\ell^b$ aporta $\ell^{\min(b,\nu)}$ a ambos
lados. Así \eqref{eq:cf-caracteres} cuenta un cociente torcido
concreto del grupo de clases.

Escribimos los divisores aditivamente. Una relación admitida consiste
en $\alpha\in F^\times$ y un divisor $\mathfrak d$ tales que
\begin{equation}
 \operatorname{div}(\alpha)=(\gamma^{-1}-a)\mathfrak d,
 \qquad
 N_a(\alpha)=\prod_{j=0}^{m_\nu-1}\gamma^j(\alpha)^{a^j}.
 \label{eq:cf-norma}
\end{equation}
El coeficiente de un primo $\mathfrak P$ en $\gamma^{-1}\mathfrak d$
es $\mathfrak d(\gamma\mathfrak P)$, lo que fija la convención
del inverso en \eqref{eq:cf-caracteres}.
La norma ponderada toma valores en $F^\times$; no se la identifica
con la norma usual hacia $L$.

El telescopado da una identidad útil:
\begin{equation}
 \operatorname{div}N_a(\alpha)
 =(1-a^{m_\nu})\gamma^{-1}\mathfrak d.
 \label{eq:cf-norma-divisor}
\end{equation}
Para comprobarla, se aplica $\sum_{j=0}^{m_\nu-1}a^j\gamma^j$
a $(\gamma^{-1}-a)\mathfrak d$. Todos los términos intermedios
se cancelan; quedan $\gamma^{-1}\mathfrak d-a^{m_\nu}\gamma^{m_\nu-1}\mathfrak d$,
y $\gamma^{m_\nu}=1$. Esta identidad explica la aparición de $a^{m_\nu}-1$
en las estimaciones de núcleos.

\subsection{El coste exacto del factor de interpolación}

Ponemos $M=2e+1$ y
\begin{equation}
 D_e=\prod_{\substack{0\le j<M\\j\ne e}}
 (\kappa^e-\kappa^j),\qquad
 B_e=2e+M(c+M).
 \label{eq:cf-interpolador}
\end{equation}
El símbolo $D_e$ designa este entero; el discriminante absoluto
sigue siendo $D$.

\begin{lema}[Valoraciones del interpolador y de la capa]
\label{lem:cf-valuaciones}
Con las condiciones sobre $c$ ya fijadas,
\begin{equation}
 \begin{aligned}
 v_\ell(\kappa^i-\kappa^j)&=c+v_\ell(i-j) &&(i\ne j),\\
 v_\ell(D_e)&=2ec+2v_\ell(e!),\\
 v_\ell(a-1)&=c+v_\ell(e),\\
 v_\ell(a^{m_\nu}-1)&=\nu+v_\ell(e).
 \end{aligned}
 \label{eq:cf-valuaciones}
\end{equation}
En particular, $D_e$ divide $\ell^{M(c+M)}$ en $\Z_\ell$.
\end{lema}

\begin{proof}
Si $b$ es primo a $\ell$, el desarrollo binomial de
$(1+\ell^c)^b-1$ tiene primer término de valoración $c$ y
los demás de valoración mayor. Para exponentes positivos,
elevar a la potencia $\ell$ aumenta esa valoración exactamente
en uno: los términos intermedios del binomio tienen mayor
valoración, también para $\ell=2$ gracias a $c\ge2$.
Escribiendo un exponente como $\ell^rb$ se obtiene
$v_\ell(\kappa^{\ell^rb}-1)=c+r$.
Los exponentes negativos se reducen a los positivos porque
$\kappa$ es una unidad. Esto prueba la primera igualdad.

En $D_e$ aparecen dos veces las distancias $1,\ldots,e$.
Sumar sus valoraciones da $2ec+2v_\ell(e!)$.
Las otras dos fórmulas aplican la primera a $e$ y
$em_\nu=e\ell^{\nu-c}$. Finalmente, cada una de las $M-1$
distancias del interpolador es menor que $M$, por lo que su
contribución es menor que $c+M$. La cota más holgada
$v_\ell(D_e)\le M(c+M)$ implica la divisibilidad en $\Z_\ell$.
\end{proof}

La valoración exacta permite conocer el coste mínimo de este paso.
Conservaremos, sin embargo, el $B_e$ de A93 en la entrada posterior
de clases y unidades: mejorar una pérdida de separación no
modifica automáticamente el enunciado aritmético que se recibe.

\subsection{Una distribución integral que aísla un peso}

El siguiente mecanismo explica por qué el interpolador aparece
en la construcción de ciclos. Para un grupo abeliano finito $A$,
su anillo de grupo $\Z[A]$ tiene base $\delta_x$ y producto
$\delta_x\delta_y=\delta_{x+y}$. Una aplicación de grupos induce
el transporte de esa base.

\begin{lema}[Aislamiento por diferencias finitas]
\label{lem:cf-distribucion}
Sean $\nu\ge c\ge1$, $Q=\ell^\nu$, $\kappa=1+\ell^c$ y
$1\le M<\ell^c$. Escribamos
$P_M(X)=\prod_{j=0}^{M-1}(X-\kappa^j)=\sum_{i=0}^Mp_iX^i$.
Para cada $t_0\in\{0,\ldots,M\}$ existe
$\eta\in\Z[(\Z/Q)^2]$ cuyo transporte por
$p_t(x,y)=x+ty$ es cero si $t\ne t_0$ en ese intervalo, y es
\begin{equation}
 (p_{t_0})_*\eta=\sum_{i=0}^Mp_i\delta_{\kappa^i}.
 \label{eq:cf-distribucion}
\end{equation}
\end{lema}

\begin{proof}
Pongamos $w=\ell^c$ y $E(f)(T)=f(T^{1+w})$ en $\Z[T]$.
Si $(T^w-1)^j$ divide $f$, entonces $(T^w-1)^{j+1}$ divide
$E(f)-\kappa^jf$. En efecto, escribimos
$f=(T^w-1)^jg$. El cociente
\[
 \frac{E(T^w-1)}{T^w-1}=\sum_{i=0}^{w}T^{wi}
 \equiv1+w\pmod{T^w-1},
\]
y $E(g)\equiv g$ módulo el mismo polinomio.
Aplicando sucesivamente $E-\kappa^j$, para $j=0,\ldots,M-1$,
al polinomio $T$, resulta
\[
 (T^w-1)^M\mid P_M(E)T=\sum_{i=0}^Mp_iT^{\kappa^i}.
\]
Evaluar $T$ en $\delta_1$ muestra que
$\lambda=\sum_i p_i\delta_{\kappa^i}$ es múltiplo de
$(\delta_w-1)^M$ en $\Z[\Z/Q]$.

Para $t\ne t_0$, la congruencia $(t_0-t)b\equiv w\pmod Q$
tiene solución porque $0<|t_0-t|\le M<\ell^c$.
Eligiendo $b\ge0$, una suma geométrica muestra que
$\delta_{t_0-t}-1$ divide $\delta_w-1$.
Hay exactamente $M$ valores de $t\ne t_0$; por tanto existe $u$
con $\lambda=\prod_{t\ne t_0}(\delta_{t_0-t}-1)u$.
No se usó que esas diferencias fueran unidades.

Sea $\iota$ la inclusión del anillo de grupo por la primera
coordenada. Tomemos
\[
 \eta=\prod_{t\ne t_0}(\delta_{(0,1)}-\delta_{(t,0)})\,
 \iota\!\left(\delta_{-\sum_{t\ne t_0}t}u\right).
\]
El transporte por un índice $t\ne t_0$ anula su factor.
Para $t_0$, cada factor se convierte en
$\delta_t(\delta_{t_0-t}-1)$. La traslación incluida cancela
el producto de los $\delta_t$, y queda $\lambda$.
\end{proof}

La desigualdad estricta $M<\ell^c$ permite resolver las
congruencias aunque alguna diferencia $t_0-t$ sea divisible por
$\ell$. La distribución tiene coeficientes enteros y soporte
finito. Realizar esos coeficientes mediante retículos de ideales,
sus mapas y sus determinantes pertenece a la entrada aritmética
de especialización que se enuncia más adelante.

\subsection{Detectar una potencia global por reducciones}

\begin{lema}[Detección por primos de grado uno]
\label{lem:cf-deteccion}
Sea $F$ un cuerpo de números con $\mu_Q\subset F$ y
$\alpha\in F^\times$. Supongamos que, fuera de un conjunto
finito de primos, su reducción es una potencia $Q$-ésima en
cada primo de $F$ de grado residual absoluto uno. Entonces
$\alpha\in(F^\times)^Q$.
\end{lema}

\begin{proof}
La extensión $K=F(\alpha^{1/Q})$ es de Galois porque $F$
contiene $\mu_Q$. Excluyamos además los primos que dividen $Q$,
el numerador o denominador de $\alpha$, y los ramificados en $K$.
En un primo restante, una raíz de $X^Q-\alpha$ en el cuerpo
residual produce todas sus raíces: las raíces de la unidad
permanecen distintas al reducir. El polinomio se descompone
completamente y Frobenius actúa trivialmente en $K/F$.

Por Chebotarev, cada clase de conjugación de
$\operatorname{Gal}(K/F)$ aparece en un conjunto de densidad
positiva. Los primos de $F$ de grado absoluto al menos dos
tienen densidad de Dirichlet cero. En efecto, si $m=[F:\Q]$,
para $s>1$ su serie de normas inversas es a lo sumo
$m\sum_p p^{-2s}$, que permanece acotada cuando $s$ tiende
a uno. Retirar esos primos y un conjunto finito de una clase
de Chebotarev deja algún primo de grado uno.
La hipótesis fuerza entonces a que cada clase de conjugación
sea la identidad. Así $K=F$ y $\alpha$ es una potencia $Q$-ésima.
\end{proof}

Se recibe el teorema clásico de Chebotarev en la forma citada
por Milne, Teorema~8.31. El grado absoluto uno tiene además
una función en la especialización: el cuerpo residual es
$\mathbf F_p$ y la imagen de $\mathcal O_L$ contiene ese cuerpo
primo, de modo que su reducción es sobreyectiva.

\subsection{De la especialización a un conjunto separado finito}

\paragraph{Entrada de ciclos y evaluación.}
Se reciben familias aditivas compatibles de evaluaciones, una
para cada cuerpo finito $k$ que contenga una raíz primitiva
de orden $Q$ y cada carácter $v:k^\times\to\Z/Q$,
normalizadas por $v$ y elegidas independientemente de $\alpha$.
Omitimos $v$ en la notación $\theta_{k,N}$.
Para cada relación \eqref{eq:cf-norma}, la construcción de A93,
(9.16), pp.~69--71, proporciona un ciclo
$y_\alpha\in C_{2e+1}(\Gamma_N,\Z/Q)$ y un
$b_\alpha\in\mathcal O_F\setminus\{0\}$.
Si $r:\mathcal O_F\to k$ es una reducción a un cuerpo finito,
su restricción a $\mathcal O_L$ es sobreyectiva y
$r(b_\alpha)\ne0$, cada familia así normalizada satisface
\begin{equation}
 \theta_{k,N}(r_*y_\alpha)
 =D_e\,v\!\left(\overline{N_a(\alpha)}\right).
 \label{eq:cf-especializacion}
\end{equation}
La existencia de la reducción del elemento fraccionario forma
parte de la afirmación. Las evaluaciones anulan fronteras,
respetan isomorfismos y suman bajo sumas directas de
representaciones sobre ciclos. Sus transportes a $\Gamma_N$
satisfacen \eqref{eq:cf-primitividad}.

El prisma y la distribución anteriores explican dos mecanismos
de esta construcción. Se conservan como entradas la realización
integral en retículos de ideales, los mapas con determinantes
prescritos, el paso a un módulo libre añadiendo un complemento
y las identidades de evaluación sobre cuerpos finitos.
Un ciclo únicamente racional no proporciona
\eqref{eq:cf-especializacion} en todos los primos de detección.

\begin{proposicion}[Separación finita]
\label{prop:cf-separacion}
Supongamos $M=2e+1<\ell^c,r_2$ y recibamos las entradas
homológicas y de especialización. Sea $S$ un conjunto finito de
elementos con relaciones \eqref{eq:cf-norma}, tales que
\begin{equation}
 \left(\frac{N_a(\alpha)}{N_a(\alpha')}\right)^{\ell^{B_e}}
 \in(F^\times)^Q
 \quad\Longrightarrow\quad\alpha=\alpha'
 \qquad(\alpha,\alpha'\in S).
 \label{eq:cf-separacion}
\end{equation}
Entonces
\begin{equation}
 \#S\le Q^{r_2}T_{4e+6,2e}(L)T_{4e+6,2e+1}(L).
 \label{eq:cf-separacion-cota}
\end{equation}
\end{proposicion}

\begin{proof}
Excluyamos los primos sobre $\ell$ y los que anulan alguno
de los finitos $b_\alpha$. En cada primo de grado absoluto uno
restante, la raíz $\zeta$ conserva orden $Q$.
El grupo $k^\times$ es cíclico; existe un carácter hacia
$\Z/Q$ cuyo núcleo es $(k^\times)^Q$.
Elijamos una familia de evaluación para cada uno y tomemos el
producto de sus transportes. Aunque haya infinitos primos,
ese producto sigue siendo un módulo anulado por $Q$ y satisface
las hipótesis del conteo homológico.

Asignemos a $\alpha$ la tupla de valores
$\ell^{2e}\theta_{k,N}(r_*y_\alpha)$.
Si dos tuplas coinciden, \eqref{eq:cf-especializacion} da
$\ell^{2e}D_e\,v(\overline{N_a(\alpha)/N_a(\alpha')})=0$
en cada reducción. La divisibilidad del
Lema~\ref{lem:cf-valuaciones} permite multiplicar por un
coeficiente de $\Z_\ell$, o por su representante módulo $Q$,
para obtener el exponente
$2e+M(c+M)=B_e$.
La reducción de la potencia en \eqref{eq:cf-separacion} queda
en el núcleo de cada carácter y es una potencia $Q$-ésima.
El Lema~\ref{lem:cf-deteccion} la hace una potencia global.
La separación implica $\alpha=\alpha'$.

La aplicación de $S$ a las tuplas es, por tanto, inyectiva.
La Proposición~\ref{prop:cf-conteo-homologico} acota su imagen
por \eqref{eq:cf-separacion-cota}. El conjunto de primos
excluidos pudo depender de $S$; la cota obtenida no depende de él.
\end{proof}

Solo para esta separación, la valoración exacta permite sustituir
$B_e$ por $2e+2ec+2v_\ell(e!)$. La comparación aritmética
siguiente se mantiene con el $B_e$ que figura en su enunciado.

\subsection{Clases, unidades y cancelación de la altura}

\paragraph{Entrada de clases y unidades.}
Escribamos $h_L=\#\operatorname{Cl}(L)$ y
$g_L=2n+\lfloor\log_2 h_L\rfloor+2$.
A93, (9.19), pp.~71--72, afirma que una cota $B$ válida para
todo conjunto finito separado como el anterior implica
\begin{equation}
 \#X_\nu\,Q^{r_2}
 \le B\,\ell^{(B_e+v_\ell(e))g_L+c+v_\ell(e)}.
 \label{eq:cf-clases-entrada}
\end{equation}
La afirmación utiliza la sucesión de divisores, las unidades,
el grupo de clases y el núcleo de la norma ponderada. El factor
$Q^{r_2}$ de la izquierda procede de esa comparación con unidades.
No se obtiene de la sola finitud del grupo de clases.

El cálculo finito que controla sus pérdidas puede verificarse
independientemente. Si un $\ell$-grupo abeliano finito $W$ cumple
$\#W[\ell]\le\ell^g$, su descomposición en cíclicos tiene a
lo sumo $g$ factores y, para $b\ge0$,
\begin{equation}
 \#W[\ell^b]\le\ell^{bg}.
 \label{eq:cf-kernel-finito}
\end{equation}
En efecto, cada factor $\Z/\ell^{a_i}$ aporta
$\ell^{\min(b,a_i)}$, y se suman a lo sumo $g$ exponentes.
Si el núcleo de un mapa desde $W$ está anulado por $\ell^b$
y su imagen tiene orden a lo sumo $B$, resulta
$\#W\le B\ell^{bg}$. Lo mismo vale para un anulador
$m\ell^b$ con $\ell\nmid m$, porque multiplicar por $m$
es un automorfismo de $W$.

En \eqref{eq:cf-clases-entrada}, identificar los grupos y probar
sus anuladores sigue siendo parte de la entrada aritmética.
Las valoraciones $c+v_\ell(e)$ y $\nu+v_\ell(e)$ de
\eqref{eq:cf-valuaciones} explican por qué ese coste conserva
$v_\ell(e)$. Una afirmación de núcleo finito sin anulador
cuantificado no produciría la desigualdad requerida.

\begin{proposicion}[Cota independiente de la altura]
\label{prop:cf-cancelacion}
Bajo las entradas anteriores y $2e+1<\ell^c,r_2$,
\begin{equation}
 \#X_\nu
 \le T_{4e+6,2e}(L)T_{4e+6,2e+1}(L)
 \ell^{(B_e+v_\ell(e))g_L+c+v_\ell(e)}
 \qquad(\nu\ge c).
 \label{eq:cf-cancelacion}
\end{equation}
Todos los términos de la derecha dependen de $L,\ell,e$,
pero no de $\nu$.
\end{proposicion}

\begin{proof}
La Proposición~\ref{prop:cf-separacion} proporciona
$B=Q^{r_2}T_{4e+6,2e}(L)T_{4e+6,2e+1}(L)$
en \eqref{eq:cf-clases-entrada}.
Como $Q^{r_2}>0$, se cancela ese mismo factor de ambos lados.
Queda \eqref{eq:cf-cancelacion}.
\end{proof}

La coincidencia de exponentes es el punto decisivo. El rango
$r_2$ de la imagen filtrada y la contribución de las unidades
se necesitan juntos. Una cota de clases que conservase
$Q^{r_2}=\ell^{\nu r_2}$ en la derecha no daría este control.

\subsection{El paso a Selmer y la extensión auxiliar}

\begin{lema}[Subgrupos anidados de orden acotado]
\label{lem:cf-limite}
Sea $A$ un grupo abeliano $\ell$-primario. Si existe una cota
finita común para $\#A[\ell^\nu]$, entonces $A$ es finito.
\end{lema}

\begin{proof}
Cada elemento de $A$ está anulado por alguna potencia de $\ell$,
así que $A=\bigcup_\nu A[\ell^\nu]$.
Los subgrupos forman una sucesión creciente y sus cardinalidades
son enteros acotados. Esos cardinales se estabilizan; una inclusión
entre dos grupos finitos del mismo orden es igualdad.
La unión es entonces uno de los subgrupos finitos.
\end{proof}

\paragraph{Entrada de comparación con Selmer.}
A93, p.~72, recibe el paso aritmético que convierte una cota
común para todos los $X_\nu$ de las capas con los datos
prescritos en
\begin{equation}
 \ell^{s_{\ell,e}(L)}
 \le T_{4e+6,2e}(L)T_{4e+6,2e+1}(L)
 \ell^{(B_e+v_\ell(e))g_L+c+v_\ell(e)}.
 \label{eq:cf-selmer-entrada}
\end{equation}
El paso conserva las condiciones locales sobre $\ell$.
El lema anterior explica el mecanismo de finitud; no identifica
los $X_\nu$ con $\operatorname{Sha}_{\ell,e}(L)[\ell^\nu]$
ni construye sus mapas de transición. Esa comparación es una
entrada adicional. La finitud para cada $L,e$ ya está disponible
por la Nota~\ref{sec:finitud-uniforme}; aquí necesitamos el
control cuantitativo que transporta la cota.

\paragraph{Entrada de extensión auxiliar.}
Las desigualdades $2e+1<\ell^c,r_2$ pueden fallar en $L$.
A93, p.~72, proporciona un cuerpo totalmente imaginario $L'$
con $n'=[L':\Q]=dn$, $D'=|\operatorname{disc}L'|$ y profundidad
ciclotómica exacta $c'$, tal que
\begin{equation}
 \begin{gathered}
 2e+1<\ell^{c'},\ r_2(L'),\qquad
 d\le(n+1)(4e+4),\qquad \ell^{c'}\le2n',\\
 D'\le D^d(\ell^{c'})^{n'},\qquad
 s_{\ell,e}(L)\le s_{\ell,e}(L').
 \end{gathered}
 \label{eq:cf-extension}
\end{equation}
Se mantiene $c'\ge1$, y $c'\ge2$ para $\ell=2$.
La última desigualdad no se atribuye a una inyectividad automática
de restricción para cualquier extensión: forma parte del cuerpo
auxiliar y de la comparación que la fuente proporciona.
El factor ciclotómico en la cota de $D'$ debe conservarse.

\subsection{Una cota suficiente para el número de clases}

\begin{proposicion}[Cota elemental a partir de Minkowski]
\label{prop:cf-clase-numero}
Para un cuerpo de números $K$ de grado $m$, discriminante
absoluto $D_K$ y número de clases $h_K$,
\begin{equation}
 \log h_K\le m^2\log2+\frac m2\log D_K.
 \label{eq:cf-clase-numero}
\end{equation}
\end{proposicion}

\begin{proof}
La cota clásica de Minkowski proporciona en cada clase un ideal
integral de norma a lo sumo $Y_0=2^m\sqrt{D_K}$;
esta forma holgada se deduce del Teorema~4.3 de Milne.
Acotemos el número de ideales de norma $a$ por el coeficiente
correspondiente de $\zeta(s)^m$.

Para un primo racional $p$, hay a lo sumo $m$ ideales primos
encima de él. Si sus grados residuales son $f_i$, los exponentes
$b_i$ de un ideal de norma $p^j$ cumplen $\sum_i f_ib_i=j$.
La aplicación $(b_i)\mapsto(f_ib_i)$, completada con ceros,
es inyectiva en las $m$-tuplas de enteros no negativos de suma
$j$. Estas últimas cuentan el coeficiente local de
$(1-p^{-s})^{-m}$. Multiplicar sobre los primos da la cota.

Sumar hasta $Y_0$ cuenta a lo sumo las $m$-tuplas de enteros
positivos cuyo producto es menor o igual que $Y_0$.
Cada coordenada es a lo sumo $\lfloor Y_0\rfloor$, por lo
que su número es a lo sumo $Y_0^m$.
Así $h_K\le Y_0^m$; tomar logaritmos prueba
\eqref{eq:cf-clase-numero}.
\end{proof}

La estimación conserva dependencia lineal en $\log D_K$.
Para la absorción final basta esa forma; optimizar la cota de
clases no sustituiría el control de la torsión homológica.

\subsection{Absorber los costes sin perder sus dependencias}

\begin{lema}[Absorción de una suma]
\label{lem:cf-absorcion}
Sean $u\ge2$, $Y\ge\log2$, $b\ge1$ y $A=u^bY$.
Entonces
\begin{equation}
 2e^A+2A\le e^{4A}\le\exp(u^{b+2}Y).
 \label{eq:cf-absorcion}
\end{equation}
\end{lema}

\begin{proof}
Tenemos $A\ge1$. La desigualdad $A\le e^A$ da
$2e^A+2A\le4e^A$. Como $\log4\le3A$,
$4e^A\le e^{4A}$. Finalmente, $u^2\ge4$.
\end{proof}

\paragraph{Entrada de torsión integral.}
La estimación de A93, p.~72, proporciona un exponente absoluto
$C_T\ge1$ tal que, para los cuerpos que intervienen y
$j=2e,2e+1$,
\begin{equation}
 \log T_{4e+6,j}(K)
 \le\exp\!\left((2e[K:\Q])^{C_T}
             \log(2|\operatorname{disc}K|)\right).
 \label{eq:cf-torsion-entrada}
\end{equation}
El orden de torsión incluye todos los primos. La construcción
del complejo integral y la estimación que producen esta entrada
son un trabajo adicional al conteo ya demostrado.

\begin{teorema}[Deducción cuantitativa de la capa fija]
\label{thm:cf-capa-fija}
Supongamos válidas las entradas de rango filtrado y estabilidad
integral, especialización de ciclos, comparación de clases y
unidades, paso cuantitativo a Selmer, extensión auxiliar y
torsión \eqref{eq:cf-torsion-entrada}, para los campos y capas
que aparecen en la construcción.
Entonces \eqref{eq:cf-objetivo} se cumple con
\begin{equation}
 C_K=\max(5C_T+9,29)+2.
 \label{eq:cf-exponente}
\end{equation}
Este valor es explícito en función del exponente recibido $C_T$.
\end{teorema}

\begin{proof}
Pongamos $u=2en$, $Y=\log(2D)$ y $Y'=\log(2D')$.
Las desigualdades de \eqref{eq:cf-extension}, con $e,n\ge1$, dan
\begin{equation}
 d\le8u,\qquad n'\le4u^2,\qquad
 2en'\le u^5,\qquad \ell^{c'}\le u^5.
 \label{eq:cf-parametros-aux}
\end{equation}
En efecto, $(n+1)(4e+4)\le16en=8u$,
$n'=dn\le8un\le4u^2$ y
$2en'\le un'\le4u^3\le u^5$.
La última cota usa $\ell^{c'}\le2n'\le8u^2\le u^5$.

El discriminante conserva el factor ramificado:
\[
 \begin{aligned}
 Y'&\le Y+d\log D+n'\log(\ell^{c'})\\
   &\le(1+8u)Y+20u^2\log u\\
   &\le(1+8u+40u^3)Y
    \le64u^3Y\le u^9Y.
 \end{aligned}
\]
Aquí se usaron $\log u\le u$, $Y>1/2$ y $u\ge2$.
Reemplazar $D'$ por $D$ directamente no sería válido.

Por la Proposición~\ref{prop:cf-clase-numero},
$\log h_{L'}\le2(n')^2Y'\le32u^{13}Y$. En consecuencia,
\begin{equation}
 \begin{aligned}
 g_{L'}&=2n'+\lfloor\log_2h_{L'}\rfloor+2\\
 &\le8u^2+2+64u^{13}Y
 \le81u^{13}Y\le u^{20}Y.
 \end{aligned}
 \label{eq:cf-generadores}
\end{equation}
Usamos $1/\log2<2$ y
$8u^2+2\le17u^2Y$, además de $u^7\ge128$.

Denotemos por $B'_e=2e+(2e+1)(c'+2e+1)$ el exponente en $L'$.
De \eqref{eq:cf-parametros-aux} se obtiene
$c'\log\ell\le5\log u$ y $\log\ell\le5\log u$.
Además, $v_\ell(e)\log\ell\le\log e\le\log u$.
Como $2e\le u$ y $2e+1\le2u$,
\[
 \begin{aligned}
 (B'_e+v_\ell(e))\log\ell
 &\le(20u^2+15u+1)\log u
 \le28u^3\le u^8,\\
 (c'+v_\ell(e))\log\ell&\le6\log u\le u^4.
 \end{aligned}
\]
El logaritmo del factor adicional de clases queda acotado por
\begin{equation}
 H'=\bigl((B'_e+v_\ell(e))g_{L'}+c'+v_\ell(e)\bigr)\log\ell
 \le u^{28}Y+u^4\le u^{29}Y.
 \label{eq:cf-resto-clases}
\end{equation}

La entrada de torsión aplicada a $L'$ y los dos grados da
$2\exp(u^{5C_T+9}Y)$ para la suma de sus logaritmos.
La comparación de Selmer \eqref{eq:cf-selmer-entrada} sobre $L'$
y \eqref{eq:cf-extension} implican
\begin{equation}
 s_{\ell,e}(L)\log\ell
 \le2\exp(u^{5C_T+9}Y)+u^{29}Y.
 \label{eq:cf-cota-sha}
\end{equation}

Falta incluir la corrección de dualidad en el cuerpo original.
La Nota~\ref{sec:finitud-uniforme} da la cota, válida también
en el primo $2$,
\[
 0\le W_\ell(L,e)-s_{\ell,e}(L)
 \le n\bigl(2+\lfloor\log_\ell n\rfloor+v_\ell(e)\bigr).
\]
Las raíces de la unidad contenidas en $L$ implican
$\ell\le n+1\le u$. Por tanto
\[
 (W_\ell-s_{\ell,e})\log\ell
 \le n(2\log\ell+\log n+\log e)
 \le2u\log u\le2u^2\le u^4Y.
\]
Tomemos $b=\max(5C_T+9,29)$ y $A=u^bY$.
Sumar esta corrección a \eqref{eq:cf-cota-sha} da
$W_\ell(L,e)\log\ell\le2e^A+2A$.
El Lema~\ref{lem:cf-absorcion} concluye
\[
 W_\ell(L,e)\log\ell
 \le\exp(u^{b+2}Y).
\]
La finitud establecida en la Nota~\ref{sec:finitud-uniforme}
permite identificar el miembro izquierdo con el logaritmo
de la cardinalidad. Esto prueba el teorema.
\end{proof}

La fórmula \eqref{eq:cf-exponente} muestra cómo absorben las
constantes; no asigna un valor numérico a $C_T$. La consecuencia
es uniforme en los parámetros indicados porque las entradas
recibidas lo son y el exponente se fija antes de variarlos.

\subsection{Qué aporta el bloque y qué entrada conserva}

Se han demostrado los pasos de cadenas, levantamiento,
conteo, separación, cancelación y absorción que enlazan las
entradas con \eqref{eq:cf-objetivo}. La tabla conserva la
procedencia de los datos que esos pasos utilizan.

\begin{center}
\small
\begin{tabularx}{\textwidth}{@{}p{0.28\textwidth}Yp{0.18\textwidth}@{}}
\toprule
Entrada recibida & Función en la deducción & Lugar en A93\\
\midrule
Rango filtrado y estabilidad integral
& Limitar los valores por fibra y conservar los órdenes de torsión.
& pp.~67--68\\
Ciclos y evaluación
& Realizar integralmente las relaciones de divisores y especializarlas.
& pp.~69--71\\
Clases y unidades
& Producir el factor $Q^{r_2}$ que cancela la altura y los anuladores.
& pp.~71--72\\
Comparación con Selmer
& Transportar la cota conservando las condiciones locales.
& p.~72\\
Extensión auxiliar
& Alcanzar el rango y las raíces requeridos con grado y discriminante controlados.
& pp.~72--73\\
Torsión en todos los primos
& Acotar los dos órdenes homológicos que permanecen tras cancelar.
& pp.~72--74\\
\bottomrule
\end{tabularx}
\end{center}

\paragraph{El alcance del antecedente de Soulé.}
Soulé, §4.1, p.~778, define $\operatorname{card}_b(A)$
eliminando de $\#A$ las partes primarias para $p\le b$.
Sus Teoremas~1--3, pp.~778--779, conservan ese recorte o
restricciones explícitas en los primos. En particular, el
Teorema~3 para el torcido $e+1$ considera
$p\ge4e+3$ y $p\ge[K:\Q]+2$.
Estos enunciados son un antecedente de las estimaciones
cuantitativas; no suministran por sí solos
\eqref{eq:cf-torsion-entrada} para toda la torsión.
El cotejo aquí se limita a sus enunciados y alcance.

\paragraph{Continuidad.}
La Nota~\ref{sec:torsion-integral} desarrolla ahora la cota integral
de torsión mediante determinantes, una cubierta de congruencia
y el descenso de índice finito. Desde el retracto de Ash y
el teorema clásico de bases acotadas obtiene $C_T=20$;
por \eqref{eq:cf-exponente}, se puede tomar $C_K=111$.
También examina el umbral impreso en el Lema~9 de Soulé
mediante un contraejemplo y una cota válida por restricción
de escalares. Las demás entradas de la tabla se desarrollan en las Notas~\ref{sec:rango-homologico-ciclos}--\ref{sec:estabilidad-integral}. En particular, el Teorema~\ref{thm:ts-lemmaK} ensambla la cota desde esas entradas, conservando sus hipótesis.
La elección de un mismo número de generadores antes de los
primos auxiliares y la comparación filtrada theta conservan
los pendientes de las Notas~\ref{sec:cohomologia},
\ref{sec:construccion-selmer} y~\ref{sec:finitud-uniforme}.

\subsection{Ejercicios con soluciones}

\begin{ejercicio}[Una obstrucción Bockstein no nula]
Sea $C_1=C_0=\Z$, con frontera multiplicación por $\ell^a$,
$a\ge1$, y los demás grupos cero. Para $Q=\ell^\nu$, $\nu\ge1$,
calcular $H_1(C/QC)$ y su Bockstein.
\end{ejercicio}

\paragraph{Solución.}
La frontera integral es inyectiva: $H_1(C)=0$ y
$H_0(C)=\Z/\ell^a$.
Los ciclos reducidos son el núcleo de multiplicar por
$\ell^a$ en $\Z/\ell^\nu$, de orden
$\ell^{\min(a,\nu)}$. Un generador es la clase de
$z=\ell^{\max(\nu-a,0)}$.
Como
\[
 \partial z
 =\ell^\nu\,\ell^{\max(a-\nu,0)},
\]
su Bockstein es la clase de $\ell^{\max(a-\nu,0)}$
en $\Z/\ell^a$. Esta clase tiene el mismo orden y genera
$H_0(C)[\ell^\nu]$.
Por ello Bockstein es un isomorfismo; ningún ciclo reducido
no nulo admite levantamiento a un ciclo integral.
La finitud del grupo reducido no elimina la obstrucción.

\begin{ejercicio}[Conteo sin generación finita]
Sea $H=\Z[1/2]\oplus\Z/9$, $J=H$ y $Q=27$.
Construir un mapa $f:H\to(\Z/27)\oplus(\Z/9)$ que alcance
la cota $Q^{\operatorname{rank}J}\#H_{\mathrm{tors}}$.
\end{ejercicio}

\paragraph{Solución.}
Como $2$ es invertible módulo $27$, la fórmula
\[
 f(a/2^b,\bar t)=(\bar a\,2^{-b},\bar t)
\]
está bien definida y es aditiva. Es sobreyectiva.
El rango de $H$ es uno y su torsión tiene orden $9$;
la imagen tiene $27\cdot9=243$ elementos, exactamente
la cota. El grupo $\Z[1/2]$ no es finitamente generado,
pero el argumento de la Proposición~\ref{prop:cf-rango-fibras}
solo necesita generar las preimágenes de un conjunto finito.

\begin{ejercicio}[Pérdida exacta y pérdida recibida]
Tomar $\ell=3$, $c=2$, $e=3$, de modo que $\kappa=10$
y $M=7<9$. Calcular $v_3(D_3)$, el exponente exacto
de filtrado y separación, y el $B_3$ de A93.
Comprobar además por qué la condición $c\ge2$ importa
en el primo $2$.
\end{ejercicio}

\paragraph{Solución.}
Las distancias a $e=3$ son $1,2,3$ dos veces; sus
valoraciones suman dos. Por tanto
\[
 v_3(D_3)=2\cdot3\cdot2+2v_3(3!)=14.
\]
El filtrado añade $2e=6$: el exponente exacto es $20$.
La versión recibida usa
$B_3=6+7(2+7)=69$.
En la comparación de clases, $v_3(e)=1$ y el factor
conservado es $3^{70g_L+3}$. Reemplazarlo por
$3^{21g_L+3}$ exigiría revisar la entrada aritmética
de núcleos, no solo el cálculo de interpolación.
Para $\ell=2$ y $c=1$,
$v_2(3^2-1)=3$, mientras que $c+v_2(2)=2$.
La fórmula general usada en el capítulo requiere
la profundidad diádica indicada.

\clearpage
\begin{ejercicio}[Un núcleo y una cancelación]
En $W=(\Z/\ell^5)\oplus(\Z/\ell^2)\oplus(\Z/\ell)$,
sea $f$ multiplicación por $\ell^2$.
Calcular núcleo e imagen y comparar con
\eqref{eq:cf-kernel-finito}.
Después, explicar qué se pierde si el conteo homológico
usa $Q^{r+1}T$ y la comparación de clases solo aporta
$Q^r$ a la izquierda.
\end{ejercicio}

\paragraph{Solución.}
El núcleo tiene orden $\ell^{2+2+1}=\ell^5$.
La imagen solo conserva $\ell^2(\Z/\ell^5)$, de orden
$\ell^3$; así $\#W=\ell^8$.
La cota con tres generadores da
$\#\ker f\le\ell^6$ y $\#W\le\ell^9$:
es válida, aunque no exacta.
En el conteo de clases,
\[
 \#X_\nu Q^r\le Q^{r+1}T\ell^A
 \quad\Longrightarrow\quad
 \#X_\nu\le Q\,T\ell^A.
\]
Permanece $Q=\ell^\nu$. La cancelación uniforme exige
controlar exactamente la potencia de la altura.
La familia $\Z/\ell^\nu$ muestra que la finitud por
separado no produce una cota común; su unión usual es
el grupo infinito $\Q_\ell/\Z_\ell$.

\begin{ejercicio}[Los primos omitidos]
Fijar un entero $b\ge2$. ¿Una cota uniforme para
$\operatorname{card}_b(A)$ acota el orden total de
un grupo abeliano finito $A$? Relacionar la respuesta
con la entrada de torsión de este capítulo.
\end{ejercicio}

\paragraph{Solución.}
Tomemos $A_m=\Z/2^m$, $m\ge1$.
Como $2\le b$, se tiene $\operatorname{card}_b(A_m)=1$
para todo $m$, mientras que $\#A_m=2^m$ no está acotado.
Incluso basta un generador para toda la familia.
Por tanto, una estimación fuera de un conjunto de primos
requiere un control adicional de sus partes primarias
antes de convertirse en una estimación del orden total.
La entrada \eqref{eq:cf-torsion-entrada} conserva precisamente
ese alcance integral.
